\begin{table}[h]
\centering
\caption{Validação estatística: densidade feminina vs masculina}
\label{tab:statistical_validation_density}
\small
\begin{tabular}{lcccl}
\hline
\textbf{Esporte} & \textbf{Densidade M} & \textbf{Densidade F} & \textbf{Razão F/M} & \textbf{Significância} \\
\hline
Atletismo 100m & 0.0152 & 0.0215 & 1.42× & $\uparrow$ \\
Natação 100m & 0.0178 & 0.0160 & 0.90× & ns \\
Basquete & 0.0243 & 0.0461 & 1.89× & $\uparrow$ \\
Futebol & 0.0133 & 0.0735 & 5.52× & $\uparrow$ \\
Judô & 0.0276 & 0.0736 & 2.66× & $\uparrow$ \\
Boxe & 0.0177 & 0.1556 & 8.78× & $\uparrow$ \\
\hline
\multicolumn{5}{l}{\textbf{Teste estatístico (Wilcoxon signed-rank pareado):}} \\
\multicolumn{5}{l}{$p = 0.031250$ *, $n = 6$ pares} \\
\multicolumn{5}{l}{Razão média: 3.53× [IC95\%: 1.57, 5.91]} \\
\hline
\end{tabular}

\vspace{0.3cm}

\footnotesize
\textit{Nota: A coluna ``Significância'' indica a direção da diferença de densidade entre os gêneros para cada esporte. $\uparrow$ = densidade feminina superior à masculina (razão F/M $>$ 1); \textbf{ns} = diferença não observada na direção esperada (razão F/M $<$ 1). O símbolo $\uparrow$ não representa um teste individual por esporte --- o teste estatístico (Wilcoxon signed-rank pareado, $p = 0.031$) é aplicado ao conjunto dos 6 pares e confirma que a direção $\uparrow$ é sistemática e estatisticamente significativa. Natação 100m constitui a única exceção, com redes masculinas ligeiramente mais densas que as femininas.}
\end{table}
